\begin{abstract}
\noindent
Activity cliffs, pairs of near-identical compounds separated by orders of magnitude in
potency, are the standard counter-example to quantitative structure-activity
relationships, and the standard remedy is a representation carrying three-dimensional
detail that a fingerprint cannot. That remedy is rarely tested against a ligand-based
baseline scored on the same data, and no study reports how large a difference its
design could have detected. We escalated four structure-based representations against
one fingerprint baseline on human acetylcholinesterase, from the geometry of a single
docked pose through an eight-member receptor ensemble and an implicit-solvent energy to
descriptors extracted from \SI{25}{\nano\second} molecular dynamics trajectories. The
ladder was scored on 271 ligands and 145 matched pairs, 77 of them constitutional
cliffs, on three framings of the endpoint, with the splits frozen beforehand and the
detection limit measured in advance. No structure-based representation was resolvably
better than the fingerprint on any framing, the design separating coefficients of
determination about $0.3$ apart. On per-ligand potency the dynamics descriptors were
resolvably worse, the gap being $-0.375$ with a 95 per cent interval of $[-0.596,
-0.184]$, while both pair endpoints gave intervals containing zero. The unexplained
variance lies far beyond what assay reproducibility accounts for and does not move with
trajectory length, which places the limitation in the descriptors rather than the data.
A replicate test found that a trajectory of this length reproduces how long a contact
is held, in 53 of 100 combinations of descriptor and analysis window, but the amplitude
of motion in only 3 of 35. What would settle the question left open is the number of
pairs and of targets, not the nanoseconds spent on each.
\end{abstract}

\noindent\textbf{Keywords.} Activity cliff; quantitative structure-activity
relationship; molecular dynamics descriptor; ensemble docking; acetylcholinesterase;
detection limit
\vspace{1em}
